% ============================================================
%  H. Høffding, "Om Realisme i Videnskab og Tro"
%  1884
%  TRANSLATION (English)
% ============================================================
\documentclass[12pt,a4paper]{article}
\usepackage[utf8]{inputenc}
\usepackage[T1]{fontenc}
\usepackage[english]{babel}
\usepackage{microtype}
\usepackage{setspace}
\usepackage[margin=1.2in]{geometry}
\usepackage{fancyhdr}
\usepackage{hyperref}

\hypersetup{
  pdftitle={On Realism in Science and Faith},
  pdfauthor={H. Høffding},
  hidelinks
}

\pagestyle{fancy}
\fancyhf{}
\rhead{Høffding, \emph{On Realism in Science and Faith} (1884)}
\cfoot{\thepage}

\setstretch{1.3}

\title{\textbf{On Realism in Science and Faith}\footnote{Essentially a
lecture delivered at the ``Studentersamfund'' on 2 September 1882.}}
\author{H.\ Høffding\\[4pt]
  \small Translated from the Danish}
\date{1884}

\begin{document}
\maketitle
\thispagestyle{fancy}

\bigskip

The word \emph{realism} is nowadays most frequently used to designate a
tendency in poetic literature. It is sometimes employed in a reproachful
sense, sometimes in a laudatory one, but rarely with a clear awareness of
what it actually expresses, or of the limits of the validity of the concept
it names. A clear account of this may perhaps only emerge once the modern
aesthetic movement has more fully and distinctly unfolded its character
and its consequences. It is not aesthetic realism that I intend to discuss
here. Rather, I shall attempt to bring out what is distinctive about the
realist tendency that, in our time, manifests itself no less clearly in
the domains of science and of general outlook on life. It follows of itself
that there must be an inner connection among the tendencies at work in
these different domains, for just as little in the spiritual world as in
the physical is there anything accidental or isolated. The human mind
cannot be partitioned in such a way that, in its imagination, it follows
paths that have absolutely nothing analogous or parallel to those it takes
in its thinking and in its faith. Only the most doctrinaire and rigid among
the opponents of modern aesthetic realism assign art and poetry such a
sharply segregated place in the life of the mind that currents from other
domains cannot---or ought not---reach them. By examining the nature and
justification of realism in one domain, we indirectly contribute to the
understanding of its appearance in others.

\section*{I.}

In the most comprehensive sense, we may define realism as the principle
of natural causes. If realism has opponents, and if it has only slowly
worked its way forward to clear self-consciousness, this must be because
the explanation of phenomena and events has been sought---and is still
being sought---by routes other than those of natural causes.

The human drive to find the causes of things has expressed itself long
before the founding of a realistic science. No period can be pointed to
in the developmental history of the human mind in which it stood entirely
passive before sense-impressions, without the slightest attempt to explain
them in some way or other. The most primitive savages have their
explanation of nature, which, seen from their standpoint, may have a
perfectly rational character. And if we could sufficiently enter into the
consciousness of an animal, we should surely find that it too interprets
the impressions it receives according to certain points of view. This
already follows from the fact that the laws governing the association of
representations prove to apply equally to animal and to human
consciousness. The history of cognition does not, therefore, begin with
a period in which single impressions are merely accumulated, to be
combined and worked over in a subsequent period. At every stage of the
development of cognition, a more or less pronounced striving manifests
itself to unite what is manifold and to fuse what is kindred.

The contrast between different periods in the history of cognition is
instead characterized by the different explanatory or interpretive
principles applied to experience. And here, above all, two great
principles stand in sharp opposition to one another. The one interprets
natural phenomena according to rules drawn from outside nature; the other
interprets nature through nature itself.

Under the first kind of natural explanation belongs all mythology. The
mythological conception of nature explains phenomena as expressions of
the will of personal beings. For it, nature is populated by more or less
human-like spirits, whose intervention is traced above all in remarkable
and significant events. From the land of fantasy and dreams, the real
world is peopled. At the primitive stage of mind, where the inhabitants
of dreamland are every bit as real as the phenomena of the actual world,
it is in itself entirely consistent and justified to lay them at the
foundation of natural explanation. In the theology of the higher folk
religions, the ultimate explanatory principle for everything is the will
of a supernatural being---likewise entirely consistent, once one has
established the reality of such a supernatural being. But the natural
explanation thereby attained does not regard nature in its own light.
Just as Archimedes wished for a point outside the earth in order to move
it, theology has sought a point outside nature in order to explain it.

Something similar holds as well of many forms of philosophical idealism.
If one finds the true explanation of a natural phenomenon in the purpose
it is assumed to serve, or in the idea it is supposed to express, one is
not explaining nature through nature itself. For whence do we know the
purpose or the idea? Experience itself does not teach us these; knowledge
of them must therefore be drawn from a source other than that through
which natural phenomena become accessible to us. We must thus possess a
special faculty for cognizing ``the Ideas,'' and our cognitive faculty
splits into two branches whose connection cannot be demonstrated. We then
have on the one side a faculty for apprehending what is given in
experience, and on the other a faculty for entering into relation with a
higher world and drawing from it the principles for understanding the
world of experience. Insofar as the principle of natural causes is
acknowledged at all, it is regarded as belonging only to the lower
cognition. True cognition comes only when we recognize the Ideas in the
world of experience. This conception, founded by Plato, the progenitor of
speculative idealism, recurs in various forms right down to the most
recent times.

Realism, with the principle of natural causes, always derives natural
phenomena from other natural phenomena, the more complex phenomena from
the simpler. When it is said that nature must be explained through itself,
it must be remembered that nature is an immeasurable world of phenomena
unfolding for us in time and in space. Every point within this world we
determine by its relation to other points. We do not go outside nature,
since we have more than enough to do in connecting the infinite number of
elements within it. Or more precisely: the expressions ``within'' and
``without'' now lose their meaning. However far we advance from cause to
effect, or however far we go back from effect to cause---no limit comes
into view. Our inquiry may come to a standstill because it lacks the
means to proceed further, as for instance the hypothesis of Kant and
Laplace takes us no further back than the rotating primordial nebula. But
wherever we begin or end, a new question always stands before us. The
principle that nature is to be explained through itself therefore sets an
infinite task. Nature as a whole, as the totality encompassing all natural
phenomena in their mutual connection, is an ideal for our cognition, which
it approaches step by step but will never be able to realize completely.

\section*{II.}

When we think through the concept of realism in this way, the absolutely
hostile opposition to idealism falls away. It is one of idealism's
fundamental thoughts that we cannot rest at the level of individual,
isolated facts, but must seek a bond that can bring coherence and unity
to bear. Idealism went wrong in assuming that we could reach this unity
and coherence by any route other than the continued working-through of
the facts of experience. In its impatience it anticipated an ideal that
can only be reached slowly and approximately; and since it did not find
this ideal fully grounded in experience, it derived it from a higher
sphere.

Realism and idealism become complete opposites only when realism
entirely abandons the maxim that genuine cognition does not consist in
accumulated experiences, but is insight into the inner connection among
experiences. But this maxim realism cannot abandon, since it is the first
and most important principle of all science. When realism turns
polemically and skeptically against idealist speculation, it can only be
because it holds that idealism has misunderstood and misused its own
maxim. Its doubt is the expression of a deeper faith: \textit{alte dubitat
qui altius credit!} Doubt and negation are not science's last word.

There are of course many popular prejudices that dissolve into nothing
under scientific criticism, many problems and questions that on closer
examination prove to owe their origin to illusory presuppositions. The
human mind is originally sanguine. It does not wait patiently for
experience's instruction, but hastens ahead with its dogmas and
hypotheses. With advancing experience there therefore also follows
advancing criticism. The development of empirical science brings with it
a damming up of those chains of representation that hitherto spread
freely in all directions like a river that has overflowed its banks. But
this negative, inhibiting activity is only one side of the matter.

The conflict between the different conceptions of nature is decided in
the same way as the conflict between dream and waking life. I attribute
reality to waking life and its contents because it is more comprehensive
and more coherent than the dream, which is always more or less
fragmentary. The greater and firmer coherence makes the lesser and looser
coherence intelligible, not the reverse. Only by giving us a more coherent
worldview than speculative idealism and theology were able to provide will
realism be able to prevail in the conflict. About particulars there will
always be room for dispute. Because a series of experiences must be
explained in accordance with realism's principle, it does not follow that
other experiences could not in themselves equally well be explained
idealistically or theologically.

Realism therefore becomes a distinctive worldview only when it ventures
beyond what is merely given, or given hitherto, and opens the prospect of
a total and thoroughgoing application of the principle of natural causes.
This is what has occurred in the great scientific hypotheses, whose series
is opened by the theory of Kant and Laplace and closes with those of
Darwin and Spencer. These hypotheses have grown naturally as consequences
of the course of development of science. In them, realism comes forward
clearly as a universal scientific principle. It is not always the case
that researchers in particular fields become conscious of this principle.
The division of labour in science has meant that the individual researcher
can with ease, and apparently also with good conscience, dig down into a
single point without troubling himself over the extent to which the
principles and methods he follows have a more general significance and
contain consequences of far-reaching importance for the overall worldview.
Many natural scientists' attitude toward Darwinism finds its explanation
here. But one who fixes his gaze on the general course of development of
the sciences and sees how every advance in understanding has been won by
applying the principle of natural causes, regards this principle as the
one called to govern our entire conception of nature. It must unfold all
its consequences; we must appropriate it; there is no way around it.

Is realism, then, proven? Not at all. And it can be proven in the strict
sense just as little as any principle can be. We arrive at a principle
when, by following a developmental course, we discover a dominant rule or
tendency that becomes more clearly visible the further the development
proceeds. Every principle is, to that extent, a hypothesis that we try to
apply to reality. We try one principle after another, until we find the
one best able to render the given intelligible. The gods of mythology and
the Ideas of Platonism would still reign in the sciences, if the advancing
experience were explained just as well by them as by the natural causes
accessible to experience itself. The validity and justification of
scientific realism will depend on whether it truly indicates the only
principle capable of bringing unity and harmony to our thinking, and on
whether it continues to receive confirmation as science advances.
Provisionally it indicates the only way in which our general worldview can
come into harmony with the spirit in which particular investigations are
conducted---which follows from the fact that it is itself nothing other
than the modern spirit of science itself.

\section*{III.}

The expression ``natural causes'' requires closer explanation. If by the
natural we understand that which can become an object of secure human
experience, then spiritual phenomena belong to nature just as much as
material ones. When a vivid emotion sets the imagination in motion, or
when by an exertion of will we recall something in our memory, we have
here a natural causal relation just as much as when bodies expand upon
the supply of heat or are compressed under strong pressure. Nature
accordingly encompasses two groups or series of phenomena, which we can
bring into relation with one another only by way of hypothesis or
speculation: the material phenomena unfolding in space, and the phenomena
of conscious life, which become immediately accessible to us only through
self-observation, and which can only be represented in spatial form in a
symbolic way. For both groups the principle of natural causes holds.
Scientific psychology seeks the inner connection among mental phenomena,
just as physics and chemistry seek the inner connection among material
phenomena. In both cases we explain nature through nature itself.

The difficulty lies, as already noted, in finding the connection between
the two domains of experience. Materialism cuts through the knot by
overlooking their difference and simply making the spiritual a form or
product of the material. But it must firmly be maintained that the
question of the relation between the spiritual and the material is an open
question within realism.

Realism can very well acknowledge that there is such a difference between
the two groups of experience that the one cannot be reduced to the other,
even though they prove to stand in intimate connection with one another.
We know of no spiritual phenomena that are not bound up with material
phenomena, but we have no common measure that could make clear to us how
a material motion can produce a thought, or conversely.

In earlier times this difficulty did not always appear so clearly. One
let the brain produce thoughts and feelings, and conversely the will set
the limbs in motion, without concern, so long as one had neither a clear
conception of the laws of material causation nor of the fundamental
distinctive character of mental phenomena. But just as the principle of
natural causes has made its way through science as a whole, so in natural
science in particular the principle has been increasingly applied that for
every material phenomenon a material cause must also be demonstrated, and
that a material cause can only have a material effect. Physiology too---the
science of organic life---works according to this principle and advances
the demand that all organic phenomena---including, therefore, everything
that takes place in our brain when we think, feel, and will---be explained
in accordance with the principle of material causes. Physiology concedes
that there is still a very long way to go before this demand can be met,
but it rejects any explanation that does not rest on this principle.
Physiology properly concerns itself with mental phenomena only insofar as
they are bound up with certain material functions. It is these material
functions that physiology is to explain; that they are accompanied by
thoughts and feelings is, strictly speaking, a matter of indifference
to it.

But if we can no longer allow thoughts and feelings to intervene in the
series of material causes, what place and significance are we to assign
to them? Their reality no one can deny; it is immediately certain to us,
more certain than anything else, for it is through thinking that we attain
all well-grounded certainty. We might rather be driven to the assumption
that the world of thought and feeling is the only real one, and that what
we call the material world is merely a product of thought that forms
itself for us unbeknownst and by virtue of an unavoidable illusion. But
even then the question would recur: how should we conceive the relation
between our consciousness-phenomena and the material phenomena we
nonetheless think of as intimately bound up with them?

There is probably no other way out than to conceive of
consciousness-phenomena as inner, psychical expressions of the same thing
that, for the outer senses, appears as a material phenomenon. The real
existence is thus only one, but we apprehend it (at certain points at
least) in a double form. Which of these forms is the original, the
constitutive one---and whether there might perhaps be deeper-lying
principles from which both are derived---these are questions we shall not
enter into here. Only this much should be emphasized: that realism does
not here foreclose more far-reaching speculations, provided they have firm
starting-points and supports in real experiences.

Realism therefore need not contest what is the genuine motive and deepest
interest of the idealistic worldview. It is quite compatible with an
effort to uphold the value and significance of spiritual life in existence.
It cannot, to be sure, accept that the spiritual is something that descends
from a world of Ideas into the material world as into a prison; but it
does not deny or diminish its reality when it conceives of it as the inner
and, for us, most significant manifestation of the same power that is at
work in outer nature. It cannot---but neither has idealism hitherto been
able to---teach us why and how mental states and activities are bound up
with material ones; it restricts itself to the modest assumption that,
since causal connection and lawfulness appear to prevail in the world, it
is presumably no accident that consciousness and brain activity are bound
up with one another, as experience shows us.

It does not know whether there is a goal toward which all existence steers,
but it finds no impossibility in itself in the assumption that this whole
great natural interconnection, with its unbreakable laws and its firmly
linked chain of causes and effects, is like a great household, a world
order in which development and dissolution alternate, but through all the
surging there nonetheless occurs an increase of what is valuable through
the arising of higher forms of life.

\section*{IV.}

Already in the last reflections we have been led from the domain of science
over into that of faith. For the most common distinguishing mark between
science and faith is that the former investigates the laws and inner
coherence of existence, while the latter asks about the ethical
significance and value of existence.

When we define the concept of faith in this way, no necessary conflict
exists between science and faith. Science grounds and explains; faith
evaluates. The common foundation of both is the reality accessible to our
experience. Its laws are sought by science, supported by the principle of
natural causes, and faith expresses the value that this law-governed
reality must have for us in virtue of our ideals. Explanation and
evaluation need not exclude one another. What stands before us as high and
glorious, as beautiful and good, may nonetheless have had its natural
developmental history. Only false values disappear when the explanation
is won. When we believe that world-development is more than a play of
atoms, that---even if it occurs through much struggle---something good and
valuable works its way forward through it, this faith does not exclude an
equally firm conviction that this result is attained only through a firmly
coherent causal sequence.

The reason we nonetheless distinguish between faith and knowledge is
partly that ideal evaluation in and of itself says nothing about the real
conditions for the arising and subsistence of what is valuable, and partly
that the conviction of the subsistence and validity of what is valuable
has its root more in feeling and will than in cognition. Only when I feel
ideal tasks and ends as the true centre of gravity of my own conduct of
life do I have occasion to conceive of natural development as progressive
realization of an ideal content. The ethical ideal I acknowledge casts its
light out over nature. As I am---thinking, feeling, and willing, with my
ideals and my capacities---I have emerged from the great natural
development. What stands fixed for me as the unconditionally true and good
is certainly conditioned by my finite nature and my vanishing position in
the great whole; but thereby it is not rendered meaningless and
accidental. It is by its fruits that one shall know the tree, and from the
fruits that nature has allowed to ripen in and for us, we can learn that
forces are at work within it that not only operate in the same direction
as our highest strivings, but have themselves from the outset led us into
these strivings.

On such a foundation, however, no theoretical system of doctrine can be
built, as speculative philosophy assumed. In its further elaboration,
poetic symbols would soon take the place of scientific concepts.

By realism in faith I mean a view of life that holds fast to the belief
in the ideal value of world-development, while at the same time being
equally deeply convinced that this ideal value is realized through the
operation of natural causes. Once this faith has been won, it is
independent of theoretical vicissitudes, for it is from the outset in
harmony with the spirit of science itself.

The customary meaning of the word faith, by contrast, excludes such a
harmony. It demands a more or less supernatural apparatus for the
realization of its ends, and sets up fixed forms and inviolable symbols.
On this rests the difference between humane or philosophical faith and
theological faith. Theological faith must, in accordance with its
principle, put forward theoretical doctrines that come into conflict with
science. It has as its object something that is supposed to be elevated
above theoretical explanation and to defy the principle of natural causes.

Under these conditions, the conflict between faith and knowledge can
never come to an end. So long as these presuppositions are maintained,
we shall not escape repetitions of the disturbing spectacle that the
history of science shows us of faith's anxiety and flight before science.
Each time that science, in experiment, proof, or hypothesis, advances to
a new domain, an outcry is raised on the part of theological faith.
Copernicus, Galileo, and Darwin, Giordano Bruno, Spinoza, and J.\ G.\
Fichte are examples of how unprejudiced observation and logical
consistency lead into conflict with a faith that has not yet learned that
evaluation and explanation do not cancel each other out. Only a faith that
can show the resignation to bow before the principle of natural causes
leads beyond this conflict. Here too it holds: \textit{alte dubitat qui
altius credit}. For it is surely the deepest faith that does not require
supernatural interventions and revelations in order to maintain the
conviction of the significance and validity of the ideal in the world.

In an entirely different sense of the word, one might also speak of a
realist tendency within theological faith in our day. Realism is
everywhere characterized by mistrust of speculation; it seeks to have
firm ground underfoot and therefore holds to what is positively given, to
the factual---or to what each variety of idealism regards as factually
given. Now when theological faith is in its own way gripped by the spirit
of realism and looks about for its facts to hold fast to, it encounters
above all its inherited formulas and confessional writings, in which it
has given to it the content sanctioned by absolute authority. It will
therefore be clear that when a realist tendency dominates an age, it will
within the domain of theological faith naturally tend to appear as
intensified faith in authority. The same tendency that produces
materialism outside the domain of theological faith produces literalism
within it. Theology has therefore also in our day abandoned every
speculative grounding and retreated to the ecclesiastically sanctioned
dogmas.

This creates a great distance between a theological thinker like Bishop
Martensen and the younger theological generation. At the end of his book
on Jakob Böhme one can read his warning words to ``the young theological
realists.'' He is surely right, for this narrow-minded clinging to
inherited dogmas is not sustainable in the long run. Nothing stands still
in the world. Even where on the surface there seems to be rest and
unchangeability, the most significant processes may be taking place in the
depths, and layers being deposited that will in time come to light. The
universally human experiences exert their quiet influence. Through the
theological-religious development since the days of the Middle Ages there
can be traced a constant tendency to reach the goal by as short and direct
a path and with as simple means as possible---the same tendency that on
scientific ground expresses itself in the rule that we must not multiply
our principles and hypotheses without compelling necessity. According to
theological faith, there are a very great many large things one must know
about in order to live as a human being and not as one who eats and drinks
because he knows he will die tomorrow. The advancing development leads
here slowly but steadily toward greater economy and modesty.
Protestantism already marks a certain reduction relative to Catholicism,
and relative to old Lutheran orthodoxy, Pietism, Rationalism, and
Grundtvigianism are all definite advances in concentration and
simplification. In this respect there is a parallel between theology and
medicine. The older medicine distinguished itself by an enormous number
of medications and a large number of heroic cures. But one gradually
manages both with fewer medications for one's bodily health and with fewer
dogmas for one's spiritual health. Where the simplification will end is
difficult to say. It will surely not end at the same point for all
individuals. But in principle the goal will surely be reached as soon as
the harmony between evaluation and explanation has properly dawned on the
general consciousness. Then the significant kernel in the higher forms of
theological faith will also be preserved, while the husks crumble away.

For what animates the higher folk religions and has given them their
mighty power over thousands of generations is ultimately the ideal impulse
that wants to see existence as something other and more than a sum of
facts, and this impulse is far nobler and more important than any of the
forms and symbols in which it has expressed itself. Humane faith can and
will preserve this impulse; indeed, its very striving is directed toward
purifying it and freeing it from all inhibiting constraints.

Only \emph{blasertude} throws away the kernel along with the husks. Such
blasertude is surely the most dangerous enemy of spiritual development in
our age, more dangerous than all intolerance and fanaticism, for in the
latter there is at least life, while blasertude is itself petrified and
petrifies everything around it. But blasertude, for which everything is
equally good so long as it is factual, naturally ends by no longer
troubling itself very much to distinguish between real and illusory facts;
it thus finally stifles the very sense for reality from which it derives
its own origin---petrifies its own source. It arises easily in times of
strong upheaval; but since it rests on a confusion of the accidental with
the essential, it will disappear, like the foam that a strong breaking of
waves produces.

\end{document}
